\chapter{Which boundary water reaches first}

This chapter writes the system that equation (1) is cut from, with the viscosity a function of
temperature and pressure, the energy balance with the heat of shear, and the boundary between
liquid and vapor as a free surface. It proves that equation (1) cannot place that boundary, and it
compares the two ways a moving body of water can reach it: by heating its liquid to the vapor
point, or by lowering its pressure to the level at which the liquid breaks. It rests on
Proposition 1 and Proposition 13, IAPWS R6-95 (2018)~\cite{src:IAPWS-R6-95-2018},
R12-08~\cite{src:IAPWS-R12-08} and R15-11~\cite{src:IAPWS-R15-11}, Caupin and Herbert's
review~\cite{src:Caupin-and-Herbert-2006}, Gruntfest and
Becker~\cite{src:Gruntfest-and-Becker-1964}, and Duraiswami~\cite{src:Duraiswami-2026}.

\section{The coupled system}

The fluid occupies a liquid region $\Omega_L(t)$ and a vapor region $\Omega_V(t)$, separated by a
surface $\Gamma(t)$. In each region the balances of the chapter on the terms equation (1) leaves out hold, with the
stress $\sigma = -p\,I + \lambda\,\theta\,I + 2\mu D$, and with every coefficient a measured
function of the state:
\[
\rho = \rho(T,p), \qquad \mu = \mu(T,p), \qquad \lambda = \lambda(T,p), \qquad k = k(T,p), \qquad
c_p = c_p(T,p).
\]
For water these are the IAPWS formulations. Without the condition $\theta = 0$, the balance of
internal energy written for the temperature is
\[
\rho\,c_p\,(\partial_t T + u\cdot\nabla T) = \beta\,T\,(\partial_t p + u\cdot\nabla p)
+ \nabla\cdot(k\nabla T) + \lambda\,\theta^2 + 2\mu\,|D|^2 ,
\]
with $\beta = -\rho^{-1}\partial\rho/\partial T$ the thermal expansion coefficient. The last two
terms are the heat of the motion, and neither is negative. These balances are the standard ones
and are not derived here.

On $\Gamma$, let $n$ be the unit normal pointing from liquid into vapor, $w$ the speed of $\Gamma$
along $n$, $\gamma$ the surface tension and $L$ the latent heat. The surface carries the balances
\[
\dot m = \rho_L\,(u_L\cdot n - w) = \rho_V\,(u_V\cdot n - w),
\]
\[
\dot m\,(u_V - u_L) - (\sigma_V - \sigma_L)\,n = \nabla_\Gamma\gamma - \gamma\,(\nabla_\Gamma\cdot n)\,n,
\]
\[
\dot m\,L = -k_L\nabla T_L\cdot n + k_V\nabla T_V\cdot n,
\]
with the temperature continuous across $\Gamma$ and equal there to the saturation temperature at
the vapor pressure. For a vapor bubble of radius $R$ at rest the second balance gives
$p_V - p_L = 2\gamma/R$, the law of Laplace.

The surface is empty until the liquid breaks somewhere. Write $p_{\mathrm{sat}}(T)$ for the vapor
pressure and $\Delta_n \ge 0$ for the tension the liquid holds below it before a vapor bubble
forms. A point of liquid stays liquid while
\[
p > p_{\mathrm{sat}}(T) - \Delta_n .
\]
The number $\Delta_n$ is a property of the liquid, of what is dissolved in it and of what it
touches, and of the volume and time over which it is held. Caupin and Herbert report that
classical nucleation theory gives $168$ MPa for water at $20\,^\circ$C, and that most methods
reach a few to a few tens of megapascals.

Set $\Gamma$ empty, the density and viscosity constant, and the energy balance aside, and the
system is equation (1), by Proposition 1. Each of those steps drops a member of
$P_{\mathrm{all}}$, and the first drops the phase.

\section{Equation (1) cannot place the boundary}

\textbf{Proposition 14.} Let $(u,p)$ solve (1) and (2) on $\R^3\times I$, with $p(\cdot,t)$
bounded for each $t\in I$, and let $p_*$ be any number. There are functions $c_1$ and $c_2$ of time
alone such that $(u,p+c_1)$ and $(u,p+c_2)$ solve (1) and (2), the set
$\{x : p(x,t) + c_1(t) \le p_*\}$ is empty for every $t$, and the set
$\{x : p(x,t) + c_2(t) \le p_*\}$ is all of $\R^3$ for every $t$.

\emph{Proof.} Take $c_1(t) = p_* + 1 - \inf_x p(x,t)$ and $c_2(t) = p_* - \sup_x p(x,t)$. Both
pairs solve (1) and (2) by Proposition 13. \hfill\qedsymbol

Whether a point of the fluid lies past the level at which the liquid breaks is therefore not a
property of a solution of (1). The same motion is liquid everywhere or past the breaking level
everywhere, according to a constant equation (1) leaves free. The coupled system fixes that
constant through $\rho(T,p)$ and the pressure the fluid is held at.

\section{The heat route}

\textbf{Proposition 15.} Let $\mu$ and $k$ be positive continuous functions of $T$, and
$F(T) = \int_{T_0}^{T} k(s)/\mu(s)\,ds$. Let the fluid fill the layer $|y| < \ell/2$ at a uniform
pressure, with velocity $(U(y),0,0)$, walls held at $T_0$ and moving at $\pm V/2$ with $V > 0$,
and steady balances of momentum and energy,
\[
(\mu(T)\,U')' = 0, \qquad (k(T)\,T')' + \mu(T)\,U'^2 = 0,
\]
with $\mu U'$ and $kT'$ continuously differentiable. If $V^2/8$ lies in the range of $F$, there
is exactly one such solution. Its temperature is largest at $y = 0$, where it is the $T_c$ with
\[
V^2 = 8\,F(T_c) = 8\int_{T_0}^{T_c}\frac{k(T)}{\mu(T)}\,dT ,
\]
which does not depend on $\ell$.

\emph{Proof.} The stress $\tau = \mu U'$ is constant, and positive since $V > 0$. The energy
balance reads $(kT')' = -\tau^2/\mu(T) < 0$. Then $kT'$ is strictly decreasing, and $T$ has one
largest value $T_c$, at the point $y_c$ where $T' = 0$. Multiplying by $kT'$ gives
$\tfrac12\big((kT')^2\big)' = -\tau^2\,F(T)'$, so that
\[
(kT')^2 = 2\tau^2\big(F(T_c) - F(T)\big).
\]
On each side of $y_c$, $T$ is monotone and $dy = k\,dT/\big(\tau\sqrt{2(F(T_c) - F(T))}\big)$. The
distance from $y_c$ to each wall is the same integral from $T_0$ to $T_c$; hence $y_c = 0$. The
velocity jump is
\[
V = \int U'\,dy = 2\int_{T_0}^{T_c}\frac{\tau}{\mu}\,
\frac{k\,dT}{\tau\sqrt{2(F(T_c) - F(T))}} = 2\int_0^{F(T_c)}\frac{dF}{\sqrt{2(F(T_c) - F)}}
= 2\sqrt{2F(T_c)} .
\]
Since $F$ is continuous and strictly increasing, $V$ fixes $T_c$. The width then fixes $\tau$
through $\ell = 2\int_{T_0}^{T_c} k\,dT/\big(\tau\sqrt{2(F(T_c)-F(T))}\big)$, an integral that
converges because $F(T_c) - F(T)$ vanishes to first order at $T_c$. The relation for $dy$ then
gives $T(y)$, and $U' = \tau/\mu(T)$ gives $U$. \hfill\qedsymbol

The heat route needs a difference of velocity, not a gradient of it: a thinner layer at the same
$V$ shears harder and conducts its heat out over a shorter distance, and its middle is no hotter.
For constant $\mu$ and $k$ the relation is $V^2 = 8k\,(T_c - T_0)/\mu$. For
$\mu = \mu_0\,e^{-a(T-T_0)}$ and constant $k$ it is
$V^2 = (8k/(a\mu_0))\,(e^{a(T_c - T_0)} - 1)$, the relation of Gruntfest and Becker. Under that law
the stress at a fixed width is proportional to $e^{-\varphi/2}\operatorname{arccosh}(e^{\varphi/2})$
with $\varphi = a(T_c - T_0)$. Its largest value is at $\varphi = 1.187$, the limit Gruntfest and
Becker give for steady flow at a fixed stress. At a fixed velocity the solution continues past it,
with the stress falling as $V$ rises.

For water at atmospheric pressure, with $T_0 = 20\,^\circ$C and $\mu(T)$ and $k(T)$ from IAPWS
R12-08 and R15-11, evaluated by programs that agree with their verification tables to six
significant figures or more, $k/\mu$ rises
from 597 to 2402 m$^2$\,s$^{-2}$\,K$^{-1}$ between $20$ and $99.9\,^\circ$C. The middle of the layer
stands at $20.012\,^\circ$C at $V = 7.63$ m/s, at $20.060\,^\circ$C at $16.9$ m/s, at
$36.7\,^\circ$C at $315$ m/s, and reaches $99.9\,^\circ$C at $956$ m/s.

\section{The tension route}

\textbf{Proposition 16.} Let $u = v(r,t)\,e_\theta$ be a swirl about an axis with no radial or
axial velocity, in a fluid of constant density $\rho$, with $p \to p_\infty$ as $r\to\infty$. Then
at each $t$
\[
p_\infty - p(0,t) = \rho\int_0^\infty \frac{v(r,t)^2}{r}\,dr .
\]
For the Lamb--Oseen vortex, $v = (\Gamma_0/2\pi r)\big(1 - e^{-r^2/r_c^2}\big)$, an exact solution
of (1) for any core radius $r_c(t)$ with $r_c^2 = r_0^2 + 4\nu t$,
\[
p_\infty - p(0,t) = \rho\Big(\frac{\Gamma_0}{2\pi r_c}\Big)^2\ln 2 = 1.702\,\rho\,v_{\max}^2 ,
\]
with $v_{\max}$ the largest swirl.

\emph{Proof.} With no radial velocity the radial component of the momentum balance is
$\partial p/\partial r = \rho v^2/r$, and integrating from $0$ to $\infty$ gives the first
statement. For the Lamb--Oseen profile set $s = r/r_c$ and then $x = s^2$:
\[
\int_0^\infty\frac{(1-e^{-s^2})^2}{s^3}\,ds = \frac12\int_0^\infty\frac{(1-e^{-x})^2}{x^2}\,dx
= \int_0^\infty\frac{e^{-x} - e^{-2x}}{x}\,dx = \ln 2,
\]
by parts and then Frullani's integral. The swirl $(1 - e^{-s^2})/s$ is largest at $s = 1.1209$,
where it is $0.63817$. Then $v_{\max} = 0.63817\,\Gamma_0/(2\pi r_c)$ and
$\ln 2/0.63817^2 = 1.702$. \hfill\qedsymbol

For water at $20\,^\circ$C, $\rho = 998.21$ kg/m$^3$ and $p_{\mathrm{sat}} = 2339$ Pa by IAPWS
R6-95. From atmospheric pressure the axis of a Lamb--Oseen vortex reaches the vapor pressure at
$v_{\max} = 7.63$ m/s, a tension of 1 MPa at $25.5$ m/s, of 30 MPa at $133$ m/s, and the $168$ MPa
of classical nucleation theory at $315$ m/s.

\section{Which comes first}

\textbf{Proposition 17.} In water at $20\,^\circ$C and atmospheric pressure, compare a layer of
Proposition 15 and a vortex of Proposition 16 at the same speed, $V = v_{\max}$. Between $7.63$
m/s, where the axis reaches the vapor pressure, and $315$ m/s, where it reaches the tension of
classical nucleation theory, the axis passes every breaking level Caupin and Herbert report, while
the middle of the layer stays below $36.7\,^\circ$C. At $7.63$ m/s the middle of the layer is at
most $0.0122$ K above its walls, and its vapor pressure is raised by less than $1.8$ Pa.

\emph{Proof.} Since $k/\mu$ is increasing on the interval, $F(T_c) \ge (k/\mu)(T_0)\,(T_c - T_0)$,
so that $T_c - T_0 \le \mu(T_0)\,V^2/(8k(T_0))$, which at $7.63$ m/s is $0.0122$ K. The vapor
pressure rises by $145$ Pa per kelvin at $20\,^\circ$C, and less than $1.8$ Pa over that interval.
The other figures are the computed values of the two previous sections. \hfill\qedsymbol

At a common speed the vortex lowers its axis pressure by $1.702\,\rho V^2$, and the layer raises the
vapor pressure at its middle by about $(dp_{\mathrm{sat}}/dT)\,\mu V^2/(8k)$. Their ratio,
$8\cdot 1.702\,\rho k/(\mu\,dp_{\mathrm{sat}}/dT)$, is $5.6\times10^4$ for water at $20\,^\circ$C.
The tension route wins by that factor, and the boundary water reaches first is the cavity. This
agrees with Duraiswami's computation for the forced construction, which puts its core at
cavitation when the swirl reaches 10 to 17 m/s, with the heat of the remaining collapse near
0.01 K.

\section{The heat route at the instant of contact}

A layer between walls held at $T_0$ loses its heat through the walls. Two sheets of one liquid
brought together with a velocity jump have no walls at the boundary between them, and the heat of
the first instant has nowhere to go but into the liquid on either side.

\textbf{Proposition 18.} Let one fluid of constant $\rho$, $\mu$, $c$ and $k$, at $T_0$, fill
$y > 0$ moving at $V/2$ along $x$ and $y < 0$ moving at $-V/2$, from $t = 0$. Write
$\nu = \mu/\rho$, $\Pr = \mu c/k$ and $\eta = y/(2\sqrt{\nu t})$. The velocity
$U = (V/2)\,\mathrm{erf}(\eta)$ and the temperature
\[
T = T_0 + \frac{V^2}{4\pi c}\,G(\eta), \qquad
G'(\eta) = -4\Pr\,e^{-\Pr\eta^2}\int_0^\eta e^{(\Pr-2)s^2}\,ds, \qquad G(\infty) = 0,
\]
solve the balances of momentum and energy. For $\Pr > 2$ the temperature on the plane $y = 0$ is
the same at every $t > 0$:
\[
T(0,t) - T_0 = \frac{V^2}{2\pi c}\,\sqrt{\frac{\Pr}{\Pr - 2}}\;
\operatorname{artanh}\sqrt{1 - \frac{2}{\Pr}} .
\]

\emph{Proof.} The velocity is the solution of Stokes's first problem, $U_t = \nu U_{yy}$, with
$U_y = (V/2)\,e^{-\eta^2}/\sqrt{\pi\nu t}$. The energy balance is
$\rho c\,T_t = k\,T_{yy} + \mu U_y^2$, since nothing varies along $x$. With
$T = T_0 + (V^2/4\pi c)\,G(\eta)$ each term carries the factor $1/t$, and dividing it out leaves
\[
\frac{G''}{4\Pr} + \frac{\eta}{2}\,G' + e^{-2\eta^2} = 0 .
\]
The integrating factor $e^{\Pr\eta^2}$ and $G'(0) = 0$, from the symmetry of the two sides, give
$G'$ as stated. Then
$G(0) = 4\Pr\int_0^\infty e^{-\Pr\eta^2}\int_0^\eta e^{(\Pr-2)s^2}\,ds\,d\eta$, and exchanging the
order of integration,
\[
G(0) = 2\sqrt{\pi\Pr}\int_0^\infty e^{(\Pr-2)s^2}\,\mathrm{erfc}\big(\sqrt{\Pr}\,s\big)\,ds
= \frac{2\sqrt{\Pr}}{\sqrt{\Pr - 2}}\,\operatorname{artanh}\sqrt{\frac{\Pr - 2}{\Pr}} ,
\]
by $\int_0^\infty e^{a s^2}\mathrm{erfc}(bs)\,ds = \operatorname{artanh}(\sqrt a/b)/\sqrt{\pi a}$
for $0 < a < b^2$. The integral converges because $\mathrm{erfc}(bs)$ falls as $e^{-b^2s^2}$.
\hfill\qedsymbol

The heat at the boundary is set by the difference in kinetic energy, $V^2/2$ per unit mass,
divided by the heat capacity, times a number fixed by $\Pr$ alone. The viscosity sets how thick
the heated layer is and not how hot it is, and neither the thickness of the sheets nor the time
since contact enters. For water at $20\,^\circ$C, $\Pr = 7.008$ and $c = 4184$
J\,kg$^{-1}$\,K$^{-1}$ by IAPWS R6-95, R12-08 and R15-11, and the boundary stands
$5.578\times10^{-5}\,V^2$ kelvin above the sheets: $0.0032$ K at $7.63$ m/s, $0.56$ K at
$100$ m/s, and $80$ K, the boiling point at atmospheric pressure, at $1198$ m/s. That speed is
$0.81$ of the speed of sound in water, where condition (2) and constant properties no longer
describe it. At $7.63$ m/s the tension route of Proposition 17 still comes first.

\section{What a blowup of (1) would have to show}

A solution of (1) whose velocity has no bound describes a liquid near its singular time only if the
liquid stays liquid there, its viscosity stays the constant (1) uses, and its density stays the
constant (1) sets. Each of these is a condition on the temperature or on the level of the pressure.

\textbf{Proposition 19.} Let $(u,p)$ be a smooth solution of (1) and (2) on $\R^3\times[0,T^*)$,
bounded with its derivatives on $\R^3\times[0,t]$ for each $t < T^*$, and let $T_b$ and $p_*$ be
numbers. Let $T$ be the bounded solution of the energy balance carried by $u$,
\[
\rho c\,(\partial_t T + u\cdot\nabla T) = k\,\Delta T + 2\mu\,|D|^2, \qquad T(\cdot,0) = T_b,
\]
with $\rho c$, $k$ and $\mu = \rho\nu$ constant. Then $T \ge T_b$ at every point and time, and with
$c_2$ from Proposition 14 the pressure $p + c_2$ is at most $p_*$ at every point and time.

\emph{Proof.} The difference $w = T - T_b$ starts at zero and solves
$\rho c\,(\partial_t w + u\cdot\nabla w) - k\,\Delta w = 2\mu|D|^2 \ge 0$, a linear parabolic
equation with bounded smooth coefficients on $\R^3\times[0,t]$. By the minimum principle for
bounded solutions of such an equation, $w \ge 0$. The statement on the pressure is
Proposition 14. \hfill\qedsymbol

Take $T_b$ to be the boiling temperature at the pressure $p_*$. The solution $(u,p)$ is then
compatible with a fluid at or above its boiling point, at a pressure at or below $p_*$, at every
point and every time: a liquid nowhere. Equation (1) is unchanged by the choice, because it holds no
temperature and fixes no level of pressure. Whether the fluid stays liquid along a solution of (1)
is therefore not a property of the solution. A theorem whose hypotheses and conclusions are
statements about $u$ and $p$ cannot decide it.

To show that a blowup of (1) describes water, a proof would have to show three things near the
singular point: that the temperature stays below boiling at the pressure there, that
$2\mu'(T)\,D\nabla T$ stays small beside $\mu\Delta u$ where $D$ has no bound (Propositions 3 and
5), and that the density holds (Proposition 4). Each needs the temperature and its gradient, and
equation (1) has neither.

\section{What this establishes and what it does not}

\textbf{Established.} The system equation (1) is cut from has a phase boundary, and equation (1)
cannot place it, because the boundary depends on a level of pressure that equation (1) leaves free.
In steady shear between walls at a fixed temperature, the heat of shear sets the middle temperature
by the velocity difference alone, for any viscosity and conductivity that depend on temperature.
In water at $20\,^\circ$C and atmospheric pressure, a turning column reaches the vapor pressure on
its axis at a speed at which a shearing layer is warmed by a hundredth of a kelvin. Two sheets of
one liquid brought together with a velocity jump $V$ heat their boundary at once, by an amount set
by $V^2$, the heat capacity and the Prandtl number, and the same at every later instant. Every
solution of (1), carried with a temperature, is compatible with a fluid past its vapor point
everywhere. Whether a blowup of (1) describes a liquid cannot be decided in (1).

\textbf{Not established.} The comparison is between two exact solutions at the same speed. It does
not show that the forced construction, or any collapsing solution, follows either profile, and a
vortex core heats at $2\mu|D|^2$ without the walls that hold the layer at $T_0$. It treats the
liquid as one substance. A boundary between two liquids that differ lowers $\Delta_n$ and the
boiling point there, and the workbook holds what that predicts in a laboratory as theory. Nothing
here shows that the coupled system is regular, or that its solutions break down.
